\documentclass[11pt]{article}

\usepackage{fontspec}
\setmainfont[Ligatures=Rare, Numbers=Proportional, Numbers=OldStyle]{Plantin MT Pro}
\newfontfamily\lasance{Lasance Initials}
\newfontfamily\canterbury{Canterbury}
\newfontfamily\vrfont{EB Garamond}
\usepackage{newunicodechar}
\newcommand{\vv}{\begin{vrfont}℣. \end{vrfont}}
\newcommand{\rr}{\begin{vrfont}℟. \end{vrfont}} 
\usepackage[autocompile]{gregoriotex}
\usepackage{multicol}
\setlength{\columnseprule}{0.5pt}
\def\columnseprulecolor{\color{black}}
\usepackage{lettrine}

\begin{document}
\gresetnabc{1}{invisible}

\begin{center}
\huge{\begin{lasance}I\end{lasance}\begin{canterbury}n\end{canterbury} \begin{lasance}N\end{lasance}\begin{canterbury}ativitate\end{canterbury} \begin{lasance}S.\end{lasance} \begin{canterbury}Joannis\end{canterbury} \begin{lasance}B\end{lasance}\begin{canterbury}aptistæ\end{canterbury}}
\end{center}

Pater noster. Ave Maria. Credo in Deum. \textit{totum secreto.}

\gregorioscore{gabc/domine}

\begin{center}
Invitatorium:
\end{center}
\gregorioscore{gabc/inv}

\begin{center}
Hymnus:
\end{center}
\gregorioscore{gabc/hymn}

\begin{center}
In I. Nocturno
\end{center}

\gregorioscore{gabc/ant1}

\begin{center}
\textit{Ps. 1}
\end{center}

\begin{multicols}{2}
\lettrine{B}{eátus} vir, qui non ábiit in consílio impiórum, † et in via peccatórum non stetit, * et in cáthedra pestiléntiæ non sedit:

Sed in lege Dómini volúntas ejus, * et in lege ejus meditábitur die ac nocte.

Et erit tamquam lignum, quod plantátum est secus decúrsus aquárum, * quod fructum suum dabit in témpore suo:

Et fólium ejus non défluet: * et ómnia quæcúmque fáciet, prosperabúntur.

Non sic ímpii, non sic: * sed tamquam pulvis, quem próicit ventus a fácie terræ.

Ídeo non resúrgent ímpii in judício: * neque peccatóres in concílio justórum.

Quóniam novit Dóminus viam justórum: * et iter impiórum períbit.

Glória Patri, et Fílio, * et Spirítui Sancto.

Sicut erat in princípio, et nunc, et semper, * et in sǽcula sæculórum. Amen.
\end{multicols}

\gregorioscore{gabc/ant1}

\gregorioscore{gabc/ant2}

\begin{center}
\textit{Ps. 2}
\end{center}

\begin{multicols}{2}
\lettrine{Q}{uare} fremuérunt gentes: * et pópuli meditáti sunt inánia?

Astitérunt reges terræ, et príncipes convenérunt in unum * advérsus Dóminum, et advérsus Christum ejus.

Dirumpámus víncula eórum: * et proiciámus a nobis jugum ipsórum.

Qui hábitat in cælis, irridébit eos: * et Dóminus subsannábit eos.

Tunc loquétur ad eos in ira sua, * et in furóre suo conturbábit eos.

Ego autem constitútus sum Rex ab eo super Sion montem sanctum ejus, * prǽdicans præcéptum ejus.

Dóminus dixit ad me: * Fílius meus es tu, ego hódie génui te.

Póstula a me, et dabo tibi gentes hereditátem tuam, * et possessiónem tuam términos terræ.

Reges eos in virga férrea, * et tamquam vas fíguli confrínges eos.

Et nunc, reges, intellégite: * erudímini, qui judicátis terram.

Servíte Dómino in timóre: * et exsultáte ei cum tremóre.

Apprehéndite disciplínam, nequándo irascátur Dóminus, * et pereátis de via justa.

Cum exárserit in brevi ira ejus: * beáti omnes qui confídunt in eo.

Glória Patri, et Fílio, * et Spirítui Sancto.

Sicut erat in princípio, et nunc, et semper, * et in sǽcula sæculórum. Amen.
\end{multicols}

\gregorioscore{gabc/ant2}

\gregorioscore{gabc/ant3}

\begin{center}
\textit{Ps. 3}
\end{center}

\begin{multicols}{2}
\lettrine{D}{ómine,} quid multiplicáti sunt qui tríbulant me? * multi insúrgunt advérsum me.

Multi dicunt ánimæ meæ: * Non est salus ipsi in Deo ejus.

Tu autem, Dómine, suscéptor meus es, * glória mea, et exáltans caput meum.

Voce mea ad Dóminum clamávi: * et exaudívit me de monte sancto suo.

Ego dormívi, et soporátus sum: * et exsurréxi, quia Dóminus suscépit me.

Non timébo míllia pópuli circumdántis me: * exsúrge, Dómine, salvum me fac, Deus meus.

Quóniam tu percussísti omnes adversántes mihi sine causa: * dentes peccatórum contrivísti.

Dómini est salus: * et super pópulum tuum benedíctio tua.

Glória Patri, et Fílio, * et Spirítui Sancto.

Sicut erat in princípio, et nunc, et semper, * et in sǽcula sæculórum. Amen.
\end{multicols}

\gregorioscore{gabc/ant3}

\bigskip

\vv Fuit homo missus a Deo.

\rr Cui nomen erat Joannes.

\bigskip

Pater noster \textit{secreto usque ad:}

\vv Et ne nos indúcas in tentatiónem:

\rr Sed líbera nos a malo.

\textit{Absol.} Exáudi, Dómine Jesu Christe, preces servórum tuórum, † et miserére nobis: * Qui cum Patre et Spíritu Sancto vivis et regnas in sǽcula sæculórum.

\rr Amen.

\bigskip

\textit{Lector:} Jube, domne, benedícere.

\textit{Bened.} Benedictióne perpétua * benedícat nos Pater ætérnus.

\rr Amen.

\begin{center}
\textit{Lectio j. Jeremias 1}
\end{center}

\begin{multicols}{2}

\begin{center}
Incípit Jeremías Prophéta.
\end{center}

\lettrine{V}{erba} Jeremíæ fílii Helcíæ de sacerdótibus qui fuérunt in Anathoth, in terra Bénjamin. Quod factum est verbum Dómini ad eum in diébus Josíæ fílii Amon regis Juda, in tertiodécimo anno regni ejus. Et factum est in diébus Jóakim fílii Josíæ regis Juda, usque ad consummatiónem undécimi anni Sedecíæ fílii Josíæ regis Juda, usque ad transmigratiónem Jerúsalem in mense quinto. Et factum est verbum Dómini ad me, dicens: Priúsquam te formárem in útero, novi te: et ántequam exíres de vulva, sanctificávi te, et Prophétam in géntibus dedi te.

\vv Tu autem, Dómine, miserére nobis.

\rr Deo grátias.
\end{multicols}

\gregorioscore{gabc/res1}

\bigskip

\textit{Lector:} Jube, domne, benedícere.

\textit{Bened.} Unigénitus Dei Fílius * nos benedícere et adjuváre dignétur.

\rr Amen.

\begin{center}
\textit{Lectio ij. Isa 26:1-6}
\end{center}

\begin{multicols}{2}
\lettrine{I}{n} die illa cantábitur cánticum istud in terra Juda: Urbs fortitúdinis nostræ Sion; salvátor ponétur in ea murus et antemurále. Aperíte portas, et ingrediátur gens justa, custódiens veritátem. Vetus error ábiit: servábis pacem; pacem, quia in te sperávimus. Sperástis in Dómino in sǽculis ætérnis; in Dómino Deo forti in perpétuum. Quia incurvábit habitántes in excélso; civitátem sublímem humiliábit: humiliábit eam usque ad terram, détrahet eam usque ad púlverem. Conculcábit eam pes, pedes páuperis, gressus egenórum.

\vv Tu autem, Dómine, miserére nobis.

\rr Deo grátias.
\end{multicols}

\gregorioscore{gabc/res2}

\bigskip

\textit{Lector:} Jube, domne, benedícere.

\textit{Bened.} Spíritus Sancti grátia * illúminet sensus et corda nostra.

\rr Amen.

\begin{center}
\textit{Lectio iij. Isa 26:7-9}
\end{center}

\begin{multicols}{2}
\lettrine{S}{émita} justi recta est, rectus callis justi ad ambulándum. Et in sémita judiciórum tuórum, Dómine, sustinúimus te: nomen tuum et memoriále tuum in desidério ánimæ. Anima mea desiderávit te in nocte, sed et spíritu meo in præcórdiis meis de mane vigilábo ad te. Cum féceris judícia tua in terra, justítiam discent habitatóres orbis.

\vv Tu autem, Dómine, miserére nobis.

\rr Deo grátias.
\end{multicols}

\gregorioscore{gabc/res3}

\begin{center}
In II. Nocturno
\end{center}

\gregorioscore{gabc/ant4}

\begin{center}
\textit{Ps. 44}
\end{center}

\begin{multicols}{2}
\lettrine{E}{ructávit} cor meum verbum bonum: * dico ego ópera mea Regi.

Lingua mea cálamus scribæ: * velóciter scribéntis.

Speciósus forma præ fíliis hóminum, diffúsa est grátia in lábiis tuis: * proptérea benedíxit te Deus in ætérnum.

Accíngere gládio tuo super femur tuum, * potentíssime.

Spécie tua et pulchritúdine tua: * inténde, próspere procéde, et regna.

Propter veritátem, et mansuetúdinem, et justítiam: * et dedúcet te mirabíliter déxtera tua.

Sagíttæ tuæ acútæ, pópuli sub te cadent: * in corda inimicórum Regis.

Sedes tua, Deus, in sǽculum sǽculi: * virga directiónis virga regni tui.

Dilexísti justítiam, et odísti iniquitátem: * proptérea unxit te Deus, Deus tuus, óleo lætítiæ præ consórtibus tuis.

Myrrha, et gutta, et cásia a vestiméntis tuis, a dómibus ebúrneis: * ex quibus delectavérunt te fíliæ regum in honóre tuo.

Ástitit regína a dextris tuis in vestítu deauráto: * circúmdata varietáte.

Audi fília, et vide, et inclína aurem tuam: * et oblivíscere pópulum tuum et domum patris tui.

Et concupíscet Rex decórem tuum: * quóniam ipse est Dóminus Deus tuus, et adorábunt eum.

Et fíliæ Tyri in munéribus * vultum tuum deprecabúntur: omnes dívites plebis.

Omnis glória ejus fíliæ Regis ab intus, * in fímbriis áureis circumamícta varietátibus.

Adducéntur Regi vírgines post eam: * próximæ ejus afferéntur tibi.

Afferéntur in lætítia et exsultatióne: * adducéntur in templum Regis.

Pro pátribus tuis nati sunt tibi fílii: * constítues eos príncipes super omnem terram.

Mémores erunt nóminis tui: * in omni generatióne et generatiónem.

Proptérea pópuli confitebúntur tibi in ætérnum: * et in sǽculum sǽculi.

Glória Patri, et Fílio, * et Spirítui Sancto.

Sicut erat in princípio, et nunc, et semper, * et in sǽcula sæculórum. Amen.
\end{multicols}

\gregorioscore{gabc/ant4}

\gregorioscore{gabc/ant5}

\begin{center}
\textit{Ps. 45}
\end{center}

\begin{multicols}{2}
\lettrine{D}{eus} noster refúgium, et virtus: * adjútor in tribulatiónibus, quæ invenérunt nos nimis.

Proptérea non timébimus dum turbábitur terra: * et transferéntur montes in cor maris.

Sonuérunt, et turbátæ sunt aquæ eórum: * conturbáti sunt montes in fortitúdine ejus.

Flúminis ímpetus lætíficat civitátem Dei: * sanctificávit tabernáculum suum Altíssimus.

Deus in médio ejus, non commovébitur: * adjuvábit eam Deus mane dilúculo.

Conturbátæ sunt gentes, et inclináta sunt regna: * dedit vocem suam, mota est terra.

Dóminus virtútum nobíscum: * suscéptor noster Deus Jacob.

Veníte, et vidéte ópera Dómini, quæ pósuit prodígia super terram: * áuferens bella usque ad finem terræ.

Arcum cónteret, et confrínget arma: * et scuta combúret igni.

Vacáte, et vidéte quóniam ego sum Deus: * exaltábor in géntibus, et exaltábor in terra.

Dóminus virtútum nobíscum: * suscéptor noster Deus Jacob.

Glória Patri, et Fílio, * et Spirítui Sancto.

Sicut erat in princípio, et nunc, et semper, * et in sǽcula sæculórum. Amen.
\end{multicols}

\gregorioscore{gabc/ant5}

\gregorioscore{gabc/ant6}

\begin{center}
\textit{Ps. 86}
\end{center}

\begin{multicols}{2}
\lettrine{F}{undaménta} ejus in móntibus sanctis: * díligit Dóminus portas Sion super ómnia tabernácula Jacob.

Gloriósa dicta sunt de te, * cívitas Dei.

Memor ero Rahab, et Babylónis * sciéntium me.

Ecce, alienígenæ, et Tyrus, et pópulus Æthíopum, * hi fuérunt illic.

Numquid Sion dicet: Homo, et homo natus est in ea: * et ipse fundávit eam Altíssimus?

Dóminus narrábit in scriptúris populórum, et príncipum: * horum, qui fuérunt in ea.

Sicut lætántium ómnium * habitátio est in te.

Glória Patri, et Fílio, * et Spirítui Sancto.

Sicut erat in princípio, et nunc, et semper, * et in sǽcula sæculórum. Amen.
\end{multicols}

\gregorioscore{gabc/ant6}

\bigskip

\vv Convérsus est furor tuus.

\rr Et consolátus es me.

\bigskip

Pater noster \textit{secreto usque ad:}

\vv Et ne nos indúcas in tentatiónem:

\rr Sed líbera nos a malo.

\textit{Absol.} Ipsíus píetas et misericórdia nos ádjuvet, * qui cum Patre et Spíritu Sancto vivit et regnat in sǽcula sæculórum.

\rr Amen.

\bigskip

\textit{Lector:} Jube, domne, benedícere.

\textit{Bened.} Deus Pater omnípotens * sit nobis propítius et clemens.

\rr Amen.

\begin{center}
\textit{Lectio iv.}
\end{center}

\begin{multicols}{2}
\begin{center}
De Sermóne Sancti Bernardi Abbátis.

\textit{Sermo 3. de Passione.}
\end{center}
\lettrine{Q}{uia} semel venímus ad Cor dulcíssimum Jesu, et bonum est nos hic esse, ne sinámus nos fácile avélli ab eo, de quo scriptum est: Recedéntes a te, in terra scribéntur. Quid autem accedéntes? Tu ipse doces nos. Tu dixísti accedéntibus ad te: Gaudéte, quia nómina vestra scripta sunt in cælo. Accedámus ergo ad te, et exsultábimus, et lætábimur in te, mémores Cordis tui. O quam bonum et quam jucúndum habitáre in Corde hoc! Quin pótius dabo ómnia, omnes cogitatiónes et afféctus mentis commutábo, jactans omne cogitátum meum in Cor Dómini Jesu, et sine fallácia illud me enútriet.

\vv Tu autem, Dómine, miserére nobis.

\rr Deo grátias.
\end{multicols}

\gregorioscore{gabc/res4}

\bigskip

\textit{Lector:} Jube, domne, benedícere.

\textit{Bened.} Christus perpétuæ * det nobis gáudia vitæ.

\rr Amen.

\bigskip

\begin{center}
\textit{Lectio v.}
\end{center}

\begin{multicols}{2}
\lettrine{A}{d} hoc templum, ad hæc Sancta sanctórum, ad hanc arcam testaménti adorábo, et laudábo nomen Dómini, dicens cum David: Invéni cor meum, ut orem Deum meum. Et ego invéni Cor regis, fratris, et amíci benígni Jesu. Et numquid non adorábo? Hoc ígitur Corde tuo, et meo, dulcíssime Jesu, invénto, orábo te Deum meum: admítte tantum in sacrárium exauditiónis tuæ preces meas, imo me totum trahe in Cor tuum. O ómnium pulchritúdine speciosíssime Jesu: ámplius lava me ab iniquitáte mea, et a peccáto meo munda me: ut purificátus per te, ad te puríssimum possim accédere, et in Corde tuo ómnibus diébus vitæ meæ mérear habitáre, et ut vidére simul, et fácere tuam váleam voluntátem!

\vv Tu autem, Dómine, miserére nobis.

\rr Deo grátias.
\end{multicols}

\gregorioscore{gabc/res5}

\bigskip

\textit{Lector:} Jube, domne, benedícere.

\textit{Bened.} Ignem sui amóris * accéndat Deus in córdibus nostris.

\rr Amen.

\bigskip

\begin{center}
\textit{Lectio vi.}
\end{center}

\begin{multicols}{2}
\lettrine{A}{d} hoc enim perforátum est latus tuum, ut nobis patéscat intróitus. Ad hoc vulnerátum est Cor tuum, ut in illo, et in te ab exterióribus perturbatiónibus absolúti habitáre possímus. Nihilóminus et proptérea vulnerátum est, ut per vulnus visíbile vulnus amóris invisíbile videámus. Quómodo hic ardor mélius osténdi potest, nisi quod non solum corpus, verum étiam ipsum Cor láncea vulnerári permísit? Carnále ergo vulnus, vulnus spirituále osténdit. Quis illud Cor tam vulnerátum non díligat? Quis tam amans non rédamet? Quis tam castum non amplectátur? Nos ígitur adhuc in córpore manéntes, quantum póssumus, amémus, redamémus, amplectámur vulnerátum nostrum, cujus ímpii agrícolæ fodérunt manus et pedes, latus et Cor: stemúsque, ut cor nostrum, durum adhuc et impœ́nitens, amóris sui vínculo constríngere, et jáculo vulneráre dignétur. Quam caritátem Christi patiéntis, et pro géneris humáni redemptióne moriéntis, atque in suæ mortis commemoratiónem instituéntis sacraméntum córporis et sánguinis sui, ut fidéles sub sanctíssimi Cordis sýmbolo devótius ac fervéntius recólant, ejusdémque fructus ubérius percípiant, Clemens décimus tértius Póntifex Máximus ejúsdem sanctíssimi Cordis festum quibúsdam peténtibus Ecclésiis celebráre permísit, ac dénique Summus Póntifex Pius nonus illud ad univérsam exténdit Ecclésiam.

\vv Tu autem, Dómine, miserére nobis.

\rr Deo grátias.
\end{multicols}

\gregorioscore{gabc/res6}

\begin{center}
In III. Nocturno
\end{center}

\gregorioscore{gabc/ant7}

\begin{center}
\textit{Ps. 96}
\end{center}

\begin{multicols}{2}
\lettrine{D}{óminus} regnávit, exsúltet terra: * læténtur ínsulæ multæ.

Nubes, et calígo in circúitu ejus: * justítia, et judícium corréctio sedis ejus.

Ignis ante ipsum præcédet, * et inflammábit in circúitu inimícos ejus.

Illuxérunt fúlgura ejus orbi terræ: * vidit, et commóta est terra.

Montes, sicut cera fluxérunt a fácie Dómini: * a fácie Dómini omnis terra.

Annuntiavérunt cæli justítiam ejus: * et vidérunt omnes pópuli glóriam ejus.

Confundántur omnes, qui adórant sculptília: * et qui gloriántur in simulácris suis.

Adoráte eum, omnes Ángeli ejus: * audívit, et lætáta est Sion.

Et exsultavérunt fíliæ Judæ, * propter judícia tua, Dómine:

Quóniam tu Dóminus Altíssimus super omnem terram: * nimis exaltátus es super omnes deos.

Qui dilígitis Dóminum, odíte malum: * custódit Dóminus ánimas sanctórum suórum, de manu peccatóris liberábit eos.

Lux orta est justo, * et rectis corde lætítia.

Lætámini, justi, in Dómino: * et confitémini memóriæ sanctificatiónis ejus.

Glória Patri, et Fílio, * et Spirítui Sancto.

Sicut erat in princípio, et nunc, et semper, * et in sǽcula sæculórum. Amen.
\end{multicols}

\gregorioscore{gabc/ant7}

\gregorioscore{gabc/ant8}

\begin{center}
\textit{Ps. 97}
\end{center}

\begin{multicols}{2}
\lettrine{C}{antáte} Dómino cánticum novum: * quia mirabília fecit.

Salvávit sibi déxtera ejus: * et brácchium sanctum ejus.

Notum fecit Dóminus salutáre suum: * in conspéctu géntium revelávit justítiam suam.

Recordátus est misericórdiæ suæ, * et veritátis suæ dómui Israël.

Vidérunt omnes términi terræ * salutáre Dei nostri.

Jubiláte Deo, omnis terra: * cantáte, et exsultáte, et psállite.

Psállite Dómino in cíthara, in cíthara et voce psalmi: * in tubis ductílibus, et voce tubæ córneæ.

Jubiláte in conspéctu regis Dómini: * moveátur mare, et plenitúdo ejus: orbis terrárum, et qui hábitant in eo.

Flúmina plaudent manu, simul montes exsultábunt a conspéctu Dómini: * quóniam venit judicáre terram.

Judicábit orbem terrárum in justítia, * et pópulos in æquitáte.

Glória Patri, et Fílio, * et Spirítui Sancto.

Sicut erat in princípio, et nunc, et semper, * et in sǽcula sæculórum. Amen.
\end{multicols}

\gregorioscore{gabc/ant8}

\gregorioscore{gabc/ant9}

\begin{center}
\textit{Ps. 107}
\end{center}

\begin{multicols}{2}
\lettrine{P}{arátum} cor meum, Deus, parátum cor meum: * cantábo, et psallam in glória mea.

Exsúrge, glória mea, exsúrge, psaltérium et cíthara: * exsúrgam dilúculo.

Confitébor tibi in pópulis, Dómine: * et psallam tibi in natiónibus.

Quia magna est super cælos misericórdia tua: * et usque ad nubes véritas tua:

Exaltáre super cælos, Deus, et super omnem terram glória tua: * ut liberéntur dilécti tui.

Salvum fac déxtera tua, et exáudi me: * Deus locútus est in sancto suo:

Exsultábo, et dívidam Síchimam, * et convállem tabernaculórum dimétiar.

Meus est Gálaad, et meus est Manásses: * et Éphraim suscéptio cápitis mei.

Juda rex meus: * Moab lebes spei meæ.

In Idumǽam exténdam calceaméntum meum: * mihi alienígenæ amíci facti sunt.

Quis dedúcet me in civitátem munítam? * quis dedúcet me usque in Idumǽam?

Nonne tu, Deus, qui repulísti nos, * et non exíbis, Deus, in virtútibus nostris?

Da nobis auxílium de tribulatióne: * quia vana salus hóminis.

In Deo faciémus virtútem: * et ipse ad níhilum dedúcet inimícos nostros.

Glória Patri, et Fílio, * et Spirítui Sancto.

Sicut erat in princípio, et nunc, et semper, * et in sǽcula sæculórum. Amen.
\end{multicols}

\gregorioscore{gabc/ant9}

\bigskip

\vv Ex Sion spécies decóris ejus.

\rr Deus noster maniféste véniet.

\bigskip

Pater noster \textit{secreto usque ad:}

\vv Et ne nos indúcas in tentatiónem:

\rr Sed líbera nos a malo.

\textit{Absol.} A vínculis peccatórum nostrórum * absólvat nos omnípotens et miséricors Dóminus.

\rr Amen.

\bigskip

\textit{Lector:} Jube, domne, benedícere.

\textit{Bened.} Evangélica léctio * sit nobis salus et protéctio.

\rr Amen.

\begin{center}
\textit{Lectio vij.}
\end{center}

\begin{multicols}{2}
\begin{center}
Léctio sancti Evangélii secúndum Joánnem.

\textit{Joan. 19:31-37}
\end{center}

\lettrine{I}{n} illo témpore: Judǽi, quóniam parascéve erat, ut non remanérent in cruce córpora sábbato (erat enim magnus dies ille sábbati), rogavérunt Pilátum, ut frangeréntur eórum crura, et tolleréntur. Et réliqua.

\end{multicols}

\begin{center}
Homilía sancti Augustíni Epíscopi

\textit{Tractatus 120 in Joannem.}
\end{center}

\begin{multicols}{2}
\lettrine{U}{nus} mílitum láncea latus ejus apéruit, et contínuo exívit sanguis et aqua. Vigilánti verbo Evangelísta usus est, ut non díceret: Latus ejus percússit, aut vulnerávit, aut quid áliud, sed apéruit; ut illic quodámmodo vitæ óstium panderétur, unde sacraménta Ecclésiæ manavérunt, sine quibus ad vitam, quæ vera vita est, non intrátur. Ille sanguis qui fusus est, in remissiónem fusus est peccatórum. Aqua illa salutáre témperat póculum: hæc et lavácrum præstat, et potum. Hoc prænuntiábat quod Noë in látere arcæ óstium fácere jussus est, quo intrárent animália, quæ non erant dilúvio peritúra, quibus præfigurabátur Ecclésia.

\vv Tu autem, Dómine, miserére nobis.

\rr Deo grátias.
\end{multicols}

\gregorioscore{gabc/res7}

\bigskip

\textit{Lector:} Jube, domne, benedícere.

\textit{Bened.} Divínum auxílium * máneat semper nobíscum.

\rr Amen.

\begin{center}
\textit{Lectio viij.}

De Homília Sancti Joannis Chrysóstomi.

\textit{Hom. 84}
\end{center}

\begin{multicols}{2}
\lettrine{V}{ides} quanti sit véritas? Per ea quæ Judǽi fáciunt prophetía implétur: ália enim ex hoc manifestáta est. Venérunt ergo mílites, et aliórum fregérunt crura, non Christi; sed ad Judæórum grátiam conciliándam, latus ejus láncea aperuérunt, et mórtuo adhuc insúltant. O péssimam voluntátem ac scelestíssimam! Noli tamen perturbári, o dilectíssime! quæ enim malo illi ánimo patrarunt, veritáti consentiébant. Prophetía hinc impléta est, inquiens: Vidébunt in quem transfixérunt. Neque hoc tantum, sed non creditúris fuit fídei arguméntum, ut Thomæ, et qui cum eo erant. Ad hæc et arcánum mystérium confirmátum est: exívit enim sanguis et aqua. Non casu et simplíciter hi fontes scaturiérunt, sed quóniam ex ambobus Ecclésia constitúta est.

\vv Tu autem, Dómine, miserére nobis.

\rr Deo grátias.
\end{multicols}

\gregorioscore{gabc/res8}

\bigskip

\textit{Lector:} Jube, domne, benedícere.

\textit{Bened.} Ad societátem cívium supernórum perdúcat nos Rex Angelórum.

\rr Amen.

\begin{center}
\textit{Lectio ix.}

\textit{Commem. S. Joannis a S. Facundo Conf.}
\end{center}

\begin{multicols}{2}
\lettrine{J}{oánnem,} Sahagúni in Hispánia, nóbili génere natum, paréntes cum diu prole caruíssent, piis opéribus et oratiónibus a Deo impetrárunt. Ab ineúnte ætáte futúræ sanctitátis indícia prǽbuit. Présbyter ordinátus, ut Deo quiétius servíret, omnes ecclesiásticos provéntus, quibus mérito auctus fúerat, sponte dimísit. Salmánticæ, cum in gravíssimum morbum incidísset, arctióris disciplínæ voto se obstrínxit, quod ut rédderet, ad cœnóbium sancti Augustíni, severióri disciplína tum máxime florens, se cóntulit; in quo admíssus, virtútibus omnis provectióres anteíbat. Salmanticénses cives, cruéntis factiónibus exagitátos, tum conciónibus tum privátis collóquiis ac vitæ sanctitáte, ad tranquillitátem redúxit, non semel a præsénti discrímine divínitus liberátus. Christum Dóminum, dum Sacrum fáceret, præséntem contuéri, ábdita cordis inspícere, ac futúra præsagíre frequens illi fuit. Dénique, mortis die prænuntiáto, sanctíssime ex hac vita migrávit, multis ante et post óbitum miráculis gloriósus. Quibus rite probátis, Alexánder octávus Sanctórum número eum adscrípsit.

\vv Tu autem, Dómine, miserére nobis.

\rr Deo grátias.
\end{multicols}

\gregorioscore{gabc/tedeum}

\begin{center}
\Large{Ad Laudes}
\end{center}

\gregorioscore{gabc/deus}

\gregorioscore{gabc/lant1}

\begin{center}
\textit{Ps. 92}
\end{center}

\begin{multicols}{2}
\lettrine{D}{óminus} regnávit, decórem indútus est: * indútus est Dóminus fortitúdinem, et præcínxit se.

Étenim firmávit orbem terræ, * qui non commovébitur.

Paráta sedes tua ex tunc: * a sǽculo tu es.

Elevavérunt flúmina, Dómine: * elevavérunt flúmina vocem suam.

Elevavérunt flúmina fluctus suos, * a vócibus aquárum multárum.

Mirábiles elatiónes maris: * mirábilis in altis Dóminus.

Testimónia tua credibília facta sunt nimis: * domum tuam decet sanctitúdo, Dómine, in longitúdinem diérum.

Glória Patri, et Fílio, * et Spirítui Sancto.

Sicut erat in princípio, et nunc, et semper, * et in sǽcula sæculórum. Amen.
\end{multicols}

\gregorioscore{gabc/lant1}

\gregorioscore{gabc/lant2}

\begin{center}
\textit{Ps. 99}
\end{center}

\begin{multicols}{2}
\lettrine{J}{ubiláte} Deo, omnis terra: * servíte Dómino in lætítia.

Introíte in conspéctu ejus, * in exsultatióne.

Scitóte quóniam Dóminus ipse est Deus: * ipse fecit nos, et non ipsi nos.

Pópulus ejus, et oves páscuæ ejus: † introíte portas ejus in confessióne, *
átria ejus in hymnis: confitémini illi.

Laudáte nomen ejus: quóniam suávis est Dóminus, † in ætérnum misericórdia ejus, * et usque in generatiónem et generatiónem véritas ejus.

Glória Patri, et Fílio, * et Spirítui Sancto.

Sicut erat in princípio, et nunc, et semper, * et in sǽcula sæculórum. Amen.
\end{multicols}

\gregorioscore{gabc/lant2}

\gregorioscore{gabc/lant3}

\begin{center}
\textit{Ps. 62}
\end{center}

\begin{multicols}{2}
\lettrine{D}{eus,} Deus meus, * ad te de luce vígilo.

\hspace{0.1cm}Sitívit in te ánima mea, * quam multiplíciter tibi caro mea.

In terra desérta, et ínvia, et inaquósa: † sic in sancto appárui tibi, * ut vidérem virtútem tuam, et glóriam tuam.

Quóniam mélior est misericórdia tua super vitas: * lábia mea laudábunt te.

Sic benedícam te in vita mea: * et in nómine tuo levábo manus meas.

Sicut ádipe et pinguédine repleátur ánima mea: * et lábiis exsultatiónis laudábit os meum.

Si memor fui tui super stratum meum, † in matutínis meditábor in te: * quia fuísti adjútor meus.

Et in velaménto alárum tuárum exsultábo, † adhǽsit ánima mea post te: * me suscépit déxtera tua.

Ipsi vero in vanum quæsiérunt ánimam meam, † introíbunt in inferióra terræ: * tradéntur in manus gládii, partes vúlpium erunt.

Rex vero lætábitur in Deo, † laudabúntur omnes qui jurant in eo: * quia obstrúctum est os loquéntium iníqua.

\begin{center}
\textit{Ps. 66}
\end{center}

\lettrine{D}{eus} misereátur nostri, et benedícat nobis: * illúminet vultum suum super nos, et misereátur nostri.

Ut cognoscámus in terra viam tuam, * in ómnibus géntibus salutáre tuum.

Confiteántur tibi pópuli, Deus: * confiteántur tibi pópuli omnes.

Læténtur et exsúltent gentes: † quóniam júdicas pópulos in æquitáte, * et gentes in terra dírigis.

Confiteántur tibi pópuli, Deus, † confiteántur tibi pópuli omnes: * terra dedit fructum suum.

Benedícat nos Deus, Deus noster, benedícat nos Deus: * et métuant eum omnes fines terræ.

Glória Patri, et Fílio, * et Spirítui Sancto.

Sicut erat in princípio, et nunc, et semper, * et in sǽcula sæculórum. Amen.
\end{multicols}

\gregorioscore{gabc/lant3}

\gregorioscore{gabc/lant4}

\begin{center}
\textit{Cant. Tri. Puer.}
\end{center}

\begin{multicols}{2}
\lettrine{B}{enedícite,} ómnia ópera Dómini, Dómino: * laudáte et superexaltáte eum in sǽcula.

Benedícite, Ángeli Dómini, Dómino: * benedícite, cæli, Dómino.

Benedícite, aquæ omnes, quæ super cælos sunt, Dómino: * benedícite, omnes virtútes Dómini, Dómino.

Benedícite, sol et luna, Dómino: * benedícite, stellæ cæli, Dómino.

Benedícite, omnis imber et ros, Dómino: * benedícite, omnes spíritus Dei, Dómino.

Benedícite, ignis et æstus, Dómino: * benedícite, frigus et æstus, Dómino.

Benedícite, rores et pruína, Dómino: * benedícite, gelu et frigus, Dómino.

Benedícite, glácies et nives, Dómino: * benedícite, noctes et dies, Dómino.

Benedícite, lux et ténebræ, Dómino: * benedícite, fúlgura et nubes, Dómino.

Benedícat terra Dóminum: * laudet et superexáltet eum in sǽcula.

Benedícite, montes et colles, Dómino: * benedícite, univérsa germinántia in terra, Dómino.

Benedícite, fontes, Dómino: * benedícite, mária et flúmina, Dómino.

Benedícite, cete, et ómnia, quæ movéntur in aquis, Dómino: * benedícite, omnes vólucres cæli, Dómino.

Benedícite, omnes béstiæ et pécora, Dómino: * benedícite, fílii hóminum, Dómino.

Benedícat Israël Dóminum: * laudet et superexáltet eum in sǽcula.

Benedícite, sacerdótes Dómini, Dómino: * benedícite, servi Dómini, Dómino.

Benedícite, spíritus, et ánimæ justórum, Dómino: * benedícite, sancti, et húmiles corde, Dómino.

Benedícite, Ananía, Azaría, Mísaël, Dómino: * laudáte et superexaltáte eum in sǽcula.

Benedicámus Patrem et Fílium cum Sancto Spíritu: * laudémus et superexaltémus eum in sǽcula.

Benedíctus es, Dómine, in firmaménto cæli: * et laudábilis, et gloriósus, et superexaltátus in sǽcula.
\end{multicols}

\gregorioscore{gabc/lant4}

\gregorioscore{gabc/lant5}

\begin{center}
\textit{Ps. 148}
\end{center}

\begin{multicols}{2}
\lettrine{L}{audáte} Dóminum de cælis: * laudáte eum in excélsis.

Laudáte eum, omnes Ángeli ejus: * laudáte eum, omnes virtútes ejus.

Laudáte eum, sol et luna: * laudáte eum, omnes stellæ et lumen.

Laudáte eum, cæli cælórum: * et aquæ omnes, quæ super cælos sunt, laudent 
nomen Dómini.

Quia ipse dixit, et facta sunt: * ipse mandávit, et creáta sunt.

Státuit ea in ætérnum, et in sǽculum sǽculi: * præcéptum pósuit, et non præteríbit.

Laudáte Dóminum de terra, * dracónes, et omnes abýssi.

Ignis, grando, nix, glácies, spíritus procellárum: * quæ fáciunt verbum ejus:

Montes, et omnes colles: * ligna fructífera, et omnes cedri.

Béstiæ, et univérsa pécora: * serpéntes, et vólucres pennátæ:

Reges terræ, et omnes pópuli: * príncipes, et omnes júdices terræ.

Júvenes, et vírgines: † senes cum junióribus laudent nomen Dómini: * quia exaltátum est nomen ejus solíus.

Conféssio ejus super cælum et terram: * et exaltávit cornu pópuli sui.

Hymnus ómnibus sanctis ejus: * fíliis Israël, pópulo appropinquánti sibi.

\begin{center}
\textit{Ps. 149}
\end{center}

\lettrine{C}{antáte} Dómino cánticum novum: * laus ejus in ecclésia sanctórum.

Lætétur Israël in eo, qui fecit eum: * et fílii Sion exsúltent in rege suo.

Laudent nomen ejus in choro: * in týmpano, et psaltério psallant ei:

Quia beneplácitum est Dómino in pópulo suo: * et exaltábit mansuétos in salútem.

Exsultábunt sancti in glória: * lætabúntur in cubílibus suis.

Exaltatiónes Dei in gútture eórum: * et gládii ancípites in mánibus eórum.

Ad faciéndam vindíctam in natiónibus: * increpatiónes in pópulis.

Ad alligándos reges eórum in compédibus: * et nóbiles eórum in mánicis férreis.

Ut fáciant in eis judícium conscríptum: * glória hæc est ómnibus sanctis ejus.

\begin{center}
\textit{Ps. 150}
\end{center}

\lettrine{L}{audáte} Dóminum in sanctis ejus: * laudáte eum in firmaménto virtútis ejus.

Laudáte eum in virtútibus ejus: * laudáte eum secúndum multitúdinem magnitúdinis ejus.

Laudáte eum in sono tubæ: * laudáte eum in psaltério, et cíthara.

Laudáte eum in týmpano, et choro: * laudáte eum in chordis, et órgano.

Laudáte eum in cýmbalis benesonántibus: † laudáte eum in cýmbalis jubilatiónis: * omnis spíritus laudet Dóminum.

Glória Patri, et Fílio, * et Spirítui Sancto.

Sicut erat in princípio, et nunc, et semper, * et in sǽcula sæculórum. Amen.
\end{multicols}

\gregorioscore{gabc/lant5}

\bigskip

\begin{center}
Cap. \hspace{0.5cm} \textit{Isa 12:2}
\end{center}

\lettrine{E}{cce} Deus Salvátor meus, fiduciáliter agam, et non timébo: quia fortitúdo mea et laus mea Dóminus, et factus est mihi in salútem. Hauriétis aquas in gáudio de fóntibus Salvatóris.

\rr Deo grátias.

\gregorioscore{gabc/lhymn}

\bigskip

\vv Vere languóres nostros ipse tulit.

\rr Et dolóres nostros ipse portávit.

\bigskip

\gregorioscore{gabc/lant6}

\begin{center}
\textit{Cant. Zach.}
\end{center}

\begin{multicols}{2}
\lettrine{B}{enedíctus} Dóminus, Deus Israël: * quia visitávit, et fecit redemptiónem plebis suæ:

Et eréxit cornu salútis nobis: * in domo David, púeri sui.

Sicut locútus est per os sanctórum, * qui a sǽculo sunt, prophetárum ejus:

Salútem ex inimícis nostris, * et de manu ómnium, qui odérunt nos.

Ad faciéndam misericórdiam cum pátribus nostris: * et memorári testaménti sui sancti.

Jusjurándum, quod jurávit ad Ábraham patrem nostrum, * datúrum se nobis:

Ut sine timóre, de manu inimicórum nostrórum liberáti, * serviámus illi.

In sanctitáte, et justítia coram ipso, * ómnibus diébus nostris.

Et tu, puer, Prophéta Altíssimi vocáberis: * præíbis enim ante fáciem Dómini, paráre vias ejus:

Ad dandam sciéntiam salútis plebi ejus: * in remissiónem peccatórum eórum:

Per víscera misericórdiæ Dei nostri: * in quibus visitávit nos, óriens ex alto:

Illumináre his, qui in ténebris, et in umbra mortis sedent: * ad dirigéndos pedes nostros in viam pacis.

Glória Patri, et Fílio, * et Spirítui Sancto.

Sicut erat in princípio, et nunc, et semper, * et in sǽcula sæculórum. Amen.
\end{multicols}

\gregorioscore{gabc/lant6}

\bigskip

\vv Dóminus vobíscum. \textit{vel:} Dómine, exáudi oratiónem meam.

\rr Et cum spíritu tuo. \textit{vel:} Et clamor meus ad te véniat.

\begin{center}
Orémus. \hspace{0.5cm} \textit{Oratio:}
\end{center}

\lettrine{C}{oncéde} quǽsumus omnípotens Deus: ut, qui in sanctíssimo dilécti Fílii tui Corde gloriántes, præcípua in nos caritátis ejus benefícia recólimus: eórum páriter et actu delectémur et fructu.
Per eúmdem Dóminum nostrum Jesum Christum Fílium tuum, qui tecum vivit et regnat in unitáte Spíritus Sancti, Deus, per ómnia sǽcula sæculórum.

\rr Amen.

\bigskip

\textit{Comm. S. Joannis a S. Facundo Conf.}

\gregorioscore{gabc/comm1}

\vv Justum dedúxit Dóminus per vias rectas.

\rr Et osténdit illi regnum Dei.

\begin{center}
Orémus. \hspace{0.5cm} \textit{Oratio:}
\end{center}

\lettrine{D}{eus,} auctor pacis et amátor caritátis, qui beátum Joánnem Confessórem tuum mirífica dissidéntes componéndi grátia decorásti: ejus méritis et intercessióne concéde; ut, in tua caritáte firmáti, nullis a te tentatiónibus separémur.

\bigskip

\textit{Comm. Ss. Basilidis, Cyrini, Naboris, et Nazarii, Mart.}

\gregorioscore{gabc/comm2}

\vv Exsultábunt Sancti in glória.

\rr Lætabúntur in cubílibus suis.

\begin{center}
Orémus \hspace{0.5cm} \textit{Oratio:}
\end{center}

\lettrine{S}{anctórum} Mártyrum tuórum Basílidis, Cyríni, Náboris atque Nazárii, quǽsumus, Dómine, natalítia nobis votíva respléndeant: et, quod illis cóntulit excelléntia sempitérna, frúctibus nostræ devotiónis accréscat. Per Dóminum nostrum Jesum Christum, Fílium tuum: qui tecum vivit et regnat in unitáte Spíritus Sancti, Deus, per ómnia sǽcula sæculórum.

\rr Amen.

\bigskip

\vv Dóminus vobíscum \textit{vel:} Dómine, exáudi oratiónem meam.

\rr Et cum spíritu tuo \textit{vel:} Et clamor meus ad te véniat.

\vv Benedicámus Dómino.

\rr Deo grátias.

\vv Fidélium ánimæ per misericórdiam Dei requiéscant in pace.

\rr Amen.

Pater noster \textit{(totum secreto.)}

\vv Dóminus det nobis suam pacem.

\rr Et vitam ætérnam. Amen.

\gregorioscore{gabc/salve}

\vv Ora pro nobis, sancta Dei Génitrix.

\rr Ut digni efficiámur promissiónibus Christi.

\begin{center}
Orémus. \hspace{0.5cm} \textit{Oratio:}
\end{center}

\lettrine{O}{mnípotens} sempitérne Deus, qui gloriósæ Vírginis Matris Maríæ corpus et ánimam, ut dignum Fílii tui habitáculum éffici mererétur, Spíritu Sancto cooperánte, præparásti: da, ut, cujus commemoratióne lætámur, ejus pia intercessióne, ab instántibus malis et a morte perpétua liberémur. Per eúmdem Christum Dóminum nostrum. Amen.

\rr Amen.

\bigskip

\vv Divínum auxílium máneat semper nobíscum.

\rr Amen.

\end{document}

